\chapter{Perimeter tangent measures and density}\label{sec:perimeter-tangent-measures}
Throughout this section $n\geq2$, and $E$ has locally finite perimeter in an open set $U\subset\mathbb R^n$.
Weak convergence of derivative measures alone does not imply convergence of their total variations. At a reduced-boundary point, polar alignment prevents the cancellation responsible for that failure. The results below complete the measure and density part of Theorem \ref{thm:perimeter-halfspace-blowup}. They expand \cite{Sim18}, Chapter 3, Theorem 4.4, equations (12)--(13), printed pp.~99--100. Countable rectifiability of the entire reduced boundary is a further theorem.

\begin{lemma}[The distributional derivative of an oriented half-space]\label{lem:halfspace-derivative-measure}
For a unit vector $e$ and $H=\{y:y\cdot e<0\}$,
\[
 D\chi_H=-e\,\mathcal H^{n-1}\llcorner e^\perp,
 \qquad |D\chi_H|=\mathcal H^{n-1}\llcorner e^\perp.
\]
These identities hold locally even though the half-space has infinite total volume and perimeter.
\end{lemma}
\begin{proof}
Rotate coordinates so $e=e_n$. For a compactly supported smooth field $X$, Fubini makes every tangential derivative integral zero. In the last coordinate the fundamental theorem of calculus gives
\[
 \int_H\operatorname{div}X\,dy
 =\int_{\mathbb R^{n-1}}X_n(y',0)\,dy'.
\]
The distributional sign convention therefore gives the first formula, using the Euclidean normalization of Hausdorff measure on a plane. Its density vector has unit norm, which gives the variation formula by the polar representation of vector measures. Orthogonal change of variables gives the formula for $e$.
\end{proof}

\begin{proposition}[Aligned polar vectors prevent loss of variation]\label{prop:aligned-vector-variation-limit}
Let $\mu_j$ be positive Radon measures on an open set $V$ with uniformly bounded mass on each relatively compact open subset. Let $|\nu_j|=1$ $\mu_j$-almost everywhere, fix a unit vector $e$, and suppose
\[
 \int_K|\nu_j-e|\,d\mu_j\longrightarrow0
\]
for every compact $K\subset V$. If the vector measures $-\nu_j\mu_j$ converge weakly locally to $-e\sigma$, where $\sigma$ is a positive Radon measure, then $\mu_j\rightharpoonup\sigma$ locally. This is the no-cancellation step implicit in \cite{Sim18}, equation (12).
\end{proposition}
\begin{proof}
For a continuous scalar test $\phi$ compactly supported in $K$,
\[
 \left|\int\phi\,d\mu_j-\int\phi\,\nu_j\cdot e\,d\mu_j\right|
 \leq\|\phi\|_\infty\int_K|\nu_j-e|\,d\mu_j\longrightarrow0.
\]
The second integral is the pairing of $-\nu_j\mu_j$ with the vector test $-\phi e$, so it tends to $\int\phi\,d\sigma$ by the assumed vector convergence. This proves the scalar weak convergence. The explicit polar hypothesis is what rules out mutually cancelling vector directions; weak convergence by itself would not suffice.
\end{proof}

\begin{theorem}[The unique perimeter tangent measure at a reduced-boundary point]\label{thm:reduced-boundary-perimeter-tangent}
Let $E$ have locally finite perimeter in $U$ and let $x\in\partial^*E$. For $F_r(y)=x+ry$, define
\[
 \mu_{x,r}=r^{1-n}(F_r^{-1})_\#|D\chi_E|
\]
on the expanding domain $F_r^{-1}(U)$. Then, locally on $\mathbb R^n$ as $r\downarrow0$,
\[
 \mu_{x,r}\rightharpoonup\mathcal H^{n-1}\llcorner\nu_E(x)^\perp.
\]
The limit is for every vanishing scale, with the fixed normalization $r^{1-n}$. This is \cite{Sim18}, equation (12), in outward-normal notation.
\end{theorem}
\begin{proof}
Write $e=\nu_E(x)$. Theorem \ref{thm:reduced-boundary-full-blowup} gives local $L^1$ convergence to $\chi_{H_x}$. The local perimeter bounds from Proposition \ref{prop:reduced-boundary-perimeter-growth}, after scaling, allow Proposition \ref{prop:bv-derivative-convergence} to identify the weak limit of vector derivatives as $D\chi_{H_x}=-e\mathcal H^{n-1}\llcorner e^\perp$ by Lemma \ref{lem:halfspace-derivative-measure}.

On each ball $B_R$, Lemmas \ref{lem:perimeter-polar-excess} and \ref{lem:perimeter-blowup-scaling} give
\[
 \int_{B_R}|\nu(x+ry)-e|\,d\mu_{x,r}(y)
 \leq\sqrt{\varepsilon(rR)}\,\mu_{x,r}(B_R)\longrightarrow0.
\]
Apply Proposition \ref{prop:aligned-vector-variation-limit} along any sequence $r_j\downarrow0$ to obtain the asserted scalar weak limit. A contradiction with a fixed continuous compactly supported test rules out failure for the full parameter $r\downarrow0$. Expanding domains cause no difficulty: every such test and a compact neighborhood of its support lie in the domain for all sufficiently small $r$.
\end{proof}

\begin{corollary}[Exact phase and perimeter densities]\label{cor:reduced-boundary-exact-densities}
At each $x\in\partial^*E$,
\[
 \lim_{r\downarrow0}\frac{|E\cap B_r(x)|}{\omega_n r^n}=\frac12,
 \qquad
 \lim_{r\downarrow0}\frac{|D\chi_E|(B_r(x))}{\omega_{n-1}r^{n-1}}=1.
\]
The nonsharp preliminary two-phase lower bound is thereby improved to an exact limit. This is the density conclusion in \cite{Sim18}, equation (13).
\end{corollary}
\begin{proof}
Apply the local $L^1$ convergence in Theorem \ref{thm:reduced-boundary-full-blowup} on $B_1$: the limit half-space occupies exactly half that ball, by reflection. Change variables to obtain the first limit. For the second, Theorem \ref{thm:reduced-boundary-perimeter-tangent} and the continuity-set consequence of weak Radon-measure convergence give $\mu_{x,r}(B_1)\to\mathcal H^{n-1}(B_1\cap e^\perp)=\omega_{n-1}$. Indeed the sphere $\partial B_1$ intersects $e^\perp$ in a set of zero $(n-1)$-dimensional measure. This argument can also be seen directly by sandwiching the ball indicator between continuous cutoffs of slightly smaller and larger concentric balls, then sending the radius difference to zero. Rescaling gives the displayed perimeter density.
\end{proof}

\begin{corollary}[Perimeter mass outside the tangent cone is negligible]\label{cor:reduced-boundary-cone-concentration}
For $x\in\partial^*E$, put $e=\nu_E(x)$. For every $0<\alpha<1$,
\[
 \lim_{r\downarrow0}r^{1-n}|D\chi_E|
 \bigl(\{z\in B_r(x):|(z-x)\cdot e|\geq\alpha|z-x|\}\bigr)=0.
\]
This is concentration of the perimeter measure about the tangent plane. It is not, by itself, a claim about every point in the topological boundary.
\end{corollary}
\begin{proof}
Under the change of variables $z=x+ry$, the set is contained in the fixed compact set
\[
 K=\{y\in\overline{B_1}:|y\cdot e|\geq\alpha|y|\}.
\]
Its intersection with $e^\perp$ is the single point zero. Because $n-1\geq1$, that point has zero $(n-1)$-dimensional measure. The upper bound for compact sets under weak convergence, applied to Theorem \ref{thm:reduced-boundary-perimeter-tangent}, therefore yields $\limsup\mu_{x,r}(K)\leq0$. This proves the assertion. Compact-set upper bounds follow from nonnegative continuous cutoffs and outer regularity, so no boundary regularity of $K$ is being assumed.
\end{proof}

\begin{proposition}[The outward normal converges with the perimeter measure]\label{prop:reduced-boundary-normal-tests}
Under the same hypotheses, for every compactly supported continuous vector field $X$ on $\mathbb R^n$,
\[
 \lim_{r\downarrow0}r^{1-n}
 \int X\bigl((z-x)/r\bigr)\cdot\nu(z)\,d|D\chi_E|(z)
 =\int_{\nu_E(x)^\perp}X(y)\cdot\nu_E(x)\,d\mathcal H^{n-1}(y).
\]
For sufficiently small $r$ the support of the original test lies in $U$. The formula holds for any measurable unit polar representative $\nu$.
\end{proposition}
\begin{proof}
Write $e=\nu_E(x)$ and use the rescaled variables. The difference between the left integral and $\int X\cdot e\,d\mu_{x,r}$ is bounded by $\|X\|_\infty\int_K|\nu(x+ry)-e|\,d\mu_{x,r}$ on a fixed compact neighborhood $K$ of its support. That error tends to zero by the excess estimate in the proof of Theorem \ref{thm:reduced-boundary-perimeter-tangent}. The remaining integral converges by that theorem. Changes of the polar field on a perimeter-null set affect none of the integrals.
\end{proof}
